\begin{exercise}[All Killing fields of Minkowski spacetime]{Cambridge Part III, 2003 exam}
Solve the Killing equation in Cartesian Minkowski coordinates and show that
the general solution consists of translations plus Lorentz transformations.
Count the independent generators in \(d\) dimensions.
\end{exercise}

\begin{exercise}[A Killing tensor and total angular momentum]{Cambridge Part III, General Relativity examples}
On a spherically symmetric spacetime, construct a rank-two Killing tensor
whose conserved quantity along geodesics is the square of total angular
momentum.  Show how it is built from the rotational Killing fields.
\end{exercise}

\begin{exercise}[Static versus stationary geometry]{Tong, Cambridge Part III Sheet 3}
Let \(\xi^\mu\) be a timelike Killing field.  Prove that the spacetime is
locally static precisely when \(\xi_{[\mu}\nabla_\nu\xi_{\rho]}=0\).  Show how
this condition permits coordinates with \(g_{0\ii}=0\).
\end{exercise}

\begin{exercise}[The de Sitter static horizon]{Cambridge Part III, 2018 exam}
Transform de Sitter spacetime to static coordinates, locate the cosmological
horizon, and compute its surface gravity from the norm of the static Killing
field.
\end{exercise}

\begin{exercise}[Null rays reach the anti-de Sitter boundary]{Cambridge Part III, 2025 exam}
In global anti-de Sitter spacetime, integrate a radial null geodesic from the
center to the conformal boundary.  Show that the coordinate travel time is
finite although the physical radial distance on a static slice is infinite.
\end{exercise}

\begin{exercise}[Symmetries of Bianchi I spacetime]{Cambridge Part III Cosmology, examples}
For
\[
\dd s^2=-\dd t^2+a_1^2(t)\dd x^2+a_2^2(t)\dd y^2+a_3^2(t)\dd z^2,
\]
identify the generic Killing fields.  Determine when additional rotations
appear and when the spatial symmetry enlarges to the Euclidean group.
\end{exercise}

\begin{exercise}[Local symmetry after a quotient]{Cambridge Part III, 2019 exam}
Identify points of Minkowski space under a discrete spatial translation to
form a flat cylinder.  Determine which Minkowski Killing fields descend to
global symmetries and explain why the quotient is locally but not globally
isometric to Minkowski space.
\end{exercise}

\begin{exercise}[Symmetric spaces and covariantly constant curvature]{Cambridge Part III, General Relativity examples}
A locally symmetric space admits a geodesic reflection about every point.
Show that this implies \(\nabla_\lambda R_{\mu\nu\rho\sigma}=0\).  Verify the
condition for maximally symmetric spaces and give a product-space example
that is locally symmetric but not maximally symmetric.
\end{exercise}
